\chapter{总结与展望}

本章对本文所做工作的全部内容进行总结，说明本文对于交通智能体系统的设计和实现有四个方面的主要贡献，即双层驱动架构、本地化模型部署与工具链集成、智能体测试方案、数字歧义问题在实验中发现的数字歧义问题。同时对目前的研究中数据集覆盖范围、工具链广度、硬件资源等方面的不足进行分析，并对今后实时数据接入、多模态感知、大模型升级、工具链扩展等方面的发展可能性进行预测。

\section{本文工作总结}

本文根据城市交通管理实际需要，对自然语言交互和时空数据智能分析进行研究，提出并实现一个基于大语言模型的交通智能体系统。按照研究目标中提出的研究系统架构设计、慢思考和快反应机制、交通预测模型集成、数据访问模块以及系统原型验证这五个方面，分别完成了从架构设计到代码实现再到实验验证的全部工作。本文主要的工作有如下几点。

第一，设计出一种以交通领域为主的智能体双层驱动结构。该架构分为慢思考层和快反应层，是研究目标中提出“慢思考”和“快反应”分开的设计思想。慢思考层是Qwen3.5-9B大模型对用户的意图进行理解并完成选择，即用户想要使用什么工具、调用什么样的指令，将抽象的自然语言需求分解为具体的任务序列，用ReAct循环进行多步推理。快反应层由工具执行节点并行调用底层数据服务，包含实时态势查询、历史趋势分析、未来流量预测这三个主要工具，对于简单的请求就直接给出结果，绕过大模型的推理过程来缩减响应时间。两层之间用LangGraph的ReAct推理循环进行编排，由条件路由边根据模型输出中是否包含工具调用指令来自动切换推理节点和工具执行节点，实现了推理和执行的分离。在工具调用的解析上，系统采用了 LangChain 原生 bind\_tools 优先、正则表达式回退的双重兼容机制，兼顾了标准 API 通道的可靠性和原生 XML 格式的灵活性。

第二，实现大模型的本地化部署和工具链集成，集成交通预测模块、数据访问模块。系统中使用了BIGCity时空预测模型做主要的预测工具，该模型用Transformer架构来建模地理、语义、时间这三种注意力机制之间的交通数据时空依赖关系，在PeMS08数据集上进行训练和验证。在模型部署上，用 vLLM 推理框架把Qwen3.5-9B原始精度模型部署到本地GPU服务器上，支持20KToken的上下文窗口，用PagedAttention来管理、前缀缓存等技术提高显存利用率和推理速度。关于 PeMS08 交通数据集，在数据访问上按照 ST 原子文件格式把多源异构数据统一成标准的.geo、.rel、.dyna文件，导入SQLite数据库使用FastAPI提供REST接口。三个工具用FastAPI微服务架构封装起来，和智能体决策层之间采用标准HTTP协议解耦通信，每一个服务都可以独立部署和扩展。

第三，设计并实施了一个包含七个场景的智能体测试方案，对系统原型进行了实验验证。在预测模型评估方面，将 BIGCity 与 HA、ARIMA、SVR、STGCN、DCRNN、GWNET 六个基线模型在 PeMS08 数据集上进行了对比，BIGCity 在 MAE、MAPE、RMSE 三项指标上均取得了最优结果。就智能体的验证而言，利用37条测试用例从意图识别、工具调度、报告质量以及异常处理这四个方面对智能体进行系统的评价。实验结果说明智能体对于所有的37个测试用例都准确地识别出用户的意图，可以适应五种不同的自然语言时间格式，在传感器查询、拥堵研判、流量预测、多工具协同等场景下都可以生成结构化的专业分析报告。智能体对于传感器不存在、时间超范围以及历史数据不足等异常情况都做了有意义的降级处理，并没有出现崩溃或者幻觉的情况。


\section{研究不足}

虽然本系统在设计理念和实验验证方面取得了初步成果，但是还存在着数据、功能以及硬件等各方面的不足，限制了该系统在更广阔的真实场景中应用。

数据集方面存在的限制。实验所用的PeMS08数据集只包含2016年7月至8月的历史数据，时间跨度只有两个月，并不能实时地进行数据的接入。两个因素一起造成这两个问题。智能体对于相对时间查询（昨天、明天等），由于时间格式解析正确，但是数据集中不存在对应日期的数据，因此查询结果为空。二是不能提供实时的交通状况查询，用户问“目前的拥堵情况怎样”时，系统只能给出数据集中对应时间段的历史数据，不能给出实时的路况信息。另外数据来源只来源于高速公路传感器网络，并且没有覆盖城市道路、交叉口、公共交通等各个交通场景，也没有涉及到交通事故、天气状况、大型活动等外部事件数据。这就意味着模型学到的交通模式只适用于高速公路的日周期变化规律，在更为复杂的城市路网中效果如何还有待检验。

工具链覆盖范围小。目前系统已经实现了态势查询、深度分析、流量预测这三个主要的功能，功能主要是数据查询和数值预测，没有涉及到更多的交通管理功能。路径规划工具可以给出出行路线推荐的问题，拥堵溯源工具可以从事故、施工、x`上游车流过载等各方面分析拥堵的原因，交通信号灯优化工具可以依据实时流量数据产生信号灯调整的建议。工具链的广度直接影响到智能体可以处理的业务范围，目前的覆盖范围还存在着很大的扩展空间。另外三个工具只能对单个传感器进行粒度上的查询，不能实现区域级数据的聚合、对比分析，例如查询整个海淀区的平均车速或者对比早高峰和晚高峰的拥堵路段数量等跨节点汇总任务也无法直接实现。

硬件资源限制。由于单张RTX 4090显存只有24GB，所以Qwen3.5-9B是目前单卡条件下最好的部署模型，但是70B级别的模型由于显存的限制不能在现有的硬件上运行。这就意味着系统对于极端复杂的交通状况的推理能力受到了一定的限制，在处理诸如需要同时考虑多个路段的拥堵发展情况、对绕行方案的效果进行评价、并且还要参照以往类似事件给出处置意见等这样带有多种要素综合考量性质的任务时，9B 模型所具备的推理深度以及精确度还存在改进之处。另外，vLLM部署时gpu-memory-utilization设为0.5，留出一半的显卡显存给其他服务使用，当并发请求过多的时候，由于显存不足而造成请求排队或者超时。如果后续要支持多用户同时访问，那么当前的单卡部署方案在并发能力上就是个瓶颈。

\section{未来展望}

就以上不足而言，未来的工作可以从系统架构、感知能力、模型能力、工具链四个方面着手，逐步把系统由一个原型验证系统发展成为可以实际部署的智能交通辅助决策平台。

系统架构上可以逐渐接入实时交通数据流，把系统由原来的基于历史数据的离线分析转变为真正的在线交通感知和决策系统。接入实时数据之后，智能体的昨天查询就会得到真实的昨天数据，“当前路况怎么样”也不再是需要猜测的问题了。实时数据流还可以实现智能体对路况的持续监测以及自动报警，即当某一段路的流量或者速度超过设定的阈值时，系统会主动把预警信息发给交通管理者。

在多模态感知方面，可以将视觉大模型（如 GPT-4o、LLaVA 等）接入系统 \upcite{ref39_GCTX202406013}。目前智能体只能处理文本、数值数据，但是交通管理者在实际工作中经常会查看监控画面来判断事故类型、拥堵情况。接入视觉模型之后，智能体就可以直接看懂交通监控画面了，判断一辆停着的车是不是因为交通拥堵造成的，是不是因为发生事故造成的，识别路面积水或者散落物等异常情况，还可以通过分析视频流中车辆的密度来辅助验证预测模型的结果。视觉信息与数值信息的融合，会给人类交通态势感知赋予更为完备的视角 \upcite{ref77_radford2021learning, ref63_JSJYGG202118006}。

在大模型升级方面，随着硬件条件改善，可以部署推理能力更强的大模型（如 DeepSeek-R1、GPT o1 系列等）\upcite{ref51_zhang2025application}。目前使用Qwen3.5-9B在大部分情况下表现很好，但是多步规划、复杂逻辑推理等还有待提高。推理模型在给出答案之前会进行长时间的自我推理，在进行复杂的交通分析任务，例如评价一个新路口的信号配时方案对缓解拥堵有多大的影响时，能够很好地完成。

在工具链的扩展上可以引入更多的交通领域专用工具。路径规划工具可以回答智能体“从A点到B点走哪条路最省时”这样一类问题。拥堵溯源工具可以分析出造成当前拥堵的原因是事故、施工还是上游车流太大。信号配时优化工具可以使得智能体根据实时的流量数据来自动给出信号灯配时调整的建议。工具链越丰富，智能体能处理的交通管理场景就越广。

多智能体协同时可以将单智能体结构扩展成多智能体系统。不同的智能体可以担任不同的角色，即实时交通监控与告警的智能体、长期交通规划与策略分析的智能体、公众出行信息服务的智能体。它们之间依靠消息传递来协同工作，比单个智能体去处理复杂的任务要高效得多，并且更具有容错性。该种多智能体结构给未来的交通智能体系统的发展提供了一种新的思路。